\documentclass[]{article}
\usepackage{graphicx}
\graphicspath{{../signatures/}}

% Imported Packages
%------------------------------------------------------------------------------
\usepackage{amssymb}
\usepackage{amstext}
\usepackage{amsthm}
\usepackage{amsmath}
\usepackage{enumerate}
\usepackage{fancyhdr}
\usepackage[margin=1in]{geometry}
\usepackage{graphicx}
%\usepackage{extarrows}
%\usepackage{setspace}
%\usepackage{xcolor}
\usepackage{color}
%------------------------------------------------------------------------------

% Header and Footer
%------------------------------------------------------------------------------
\pagestyle{plain}  
\renewcommand\headrulewidth{0.4pt}                                      
\renewcommand\footrulewidth{0.4pt}                                    
%------------------------------------------------------------------------------

% Title Details
%------------------------------------------------------------------------------
\title{Deliverable \#1 Template : Software Requirement Specification (SRS)}
\author{SE 3A04: Software Design II -- Large System Design}
\date{18 February 2024}
                            
%------------------------------------------------------------------------------

% Document
%------------------------------------------------------------------------------
\begin{document}

\maketitle	
\noindent{\bf Tutorial Number:} T01\\
{\bf Group Number:} G07 \\
{\bf Group Members:} 
\begin{itemize}
	\item Awurama Nyarko
	\item Chelsea Maramot
 	\item Harrison Chiu
  	\item Khushi Bhojane
   	\item Sumanya Gulati
\end{itemize}

\section*{IMPORTANT NOTES}
\begin{itemize}
	\item Be sure to include all sections of the template in your document regardless whether you have something to write for each or not
	\begin{itemize}
		\item If you do not have anything to write in a section, indicate this by the \emph{N/A}, \emph{void}, \emph{none}, etc.
	\end{itemize}
	\item Uniquely number each of your requirements for easy identification and cross-referencing
	\item Highlight terms that are defined in Section~1.3 (\textbf{Definitions, Acronyms, and Abbreviations}) with \textbf{bold}, \emph{italic} or \underline{underline}
	\item For Deliverable 1, please highlight, in some fashion, all (you may have more than one) creative and innovative features. Your creative and innovative features will generally be described in Section~2.2 (\textbf{Product Functions}), but it will depend on the type of creative or innovative features you are including.
\end{itemize}

\newpage
\section{Introduction}
\label{sec:introduction}
% Begin Section
This Software Requirements Specification (SRS) will outline the requirements and guidelines for the development of SecureChat. SecureChat is a secure communication application designed specifically for the purpose of organizational usage. This document will be discussing the applications purpose, scope, user characteristics, and an overview of the product’s overall functionalities. SecureChat will be employing robust cryptographic measures, as well as secure authentication protocols to protect sensitive information/ data and ensure secure internal communications.


\subsection{Purpose}
\label{sub:purpose}
% Begin SubSection
This SRS focuses on functional requirements, non-functional requirements, user interactions, and security protocols for the SecureChat application.\\
This document is written for the intention of usage by a broad audience including project managers, software developers, security analysts, and stakeholders. No prerequisite knowledge or readings is expected or required.
% End SubSection

\subsection{Scope}
\label{sub:scope}
% Begin SubSection
The SecureChat application is the primary software product to be developed. It is designed as an encrypted messaging app for deployment on company-issued Android devices.\\
SecureChat will integrate a Key Distribution Centre (KDC) to manage cryptographic keys and employ an authentication protocol to validate user identities. It will utilize symmetric-key encryption to ensure secure communications within the app. The application is designed to facilitate secure text messaging, file sharing, and guarantee the integrity and confidentiality of chat logs. Additional functionalities include dynamic key generation and rotation, enhancing the app's resilience against potential security threats.\\
The objective of SecureChat is to establish a secure, efficient communication channel within an organization, mitigating risks associated with corporate espionage. By utilizing advanced cryptographic techniques, SecureChat aims to protect sensitive information while maintaining user-friendly operations. Additionally, it seeks to foster trust among employees regarding the organization's commitment to security and privacy. Through these measures, SecureChat aims to offer a robust solution for secure internal communications, thereby protecting against data breaches and unauthorized access to sensitive information.
% End SubSection

\subsection{Definitions, Acronyms, and Abbreviations}
\label{sub:definitions_acronyms_and_abbreviations}
% Begin SubSection
\begin{itemize}
    \item \textbf{KDC (Key Distribution Center):} A central authority in cryptographic systems that is responsible for distributing secret keys to parties wishing to communicate securely.
    \item \textbf{Multimedia Files:} Digital files that incorporate content in multiple forms such as text, audio, images, animations, video, and interactive content.
    \item \textbf{PIPEDA (Personal Information Protection and Electronic Documents Act):} A Canadian federal privacy law regulating the collection, use, and disclosure of personal information in the course of commercial activities.
    \item \textbf{WCAG (Web Content Accessibility Guidelines):} A set of guidelines developed through the W3C process, aimed at making web content more accessible to people with disabilities.
\end{itemize}
% End SubSection

\subsection{References}
\label{sub:references}
% Begin SubSection
\begin{itemize}
    \item [1] "Personal Information Protection and Electronic Documents Act (PIPEDA)," Canada, 2000. [Online]. Available: https://www.priv.gc.ca/en/privacy-topics/privacy-laws-in-canada/the-personal-information-protection-and-electronic-documents-act-pipeda/. [Accessed: 15 Feb. 2024].
\end{itemize}
% End SubSection


\subsection{Overview}
\label{sub:overview}
% Begin SubSection
Following this introduction, Section 2 provides an in-depth description of the product, outlining its perspective, functionalities, and the general factors influencing its design and requirements. Section 3 presents the Use Case Diagram, offering a visual representation of the system's operational scenarios. Section 4 will go into the Highlights of Functional Requirements, detailing the actions and processes SecureChat is designed to perform. Categorizing them with business events and viewpoints. Section 5 will detail the Non-Functional Requirements, including the system's performance, usability, and security standards. Comprehensively these sections construct a framework for the development of SecureChat, ensuring there is clarity and a sense of direction throughout the project lifecycle.
% End SubSection

% End Section

\section{Overall Product Description}
\label{sec:overall_description}
% Begin Section

\begin{itemize}
	\item This section should describe the general factors that affect the product and its requirements. 
	\item It does not state specific requirements.
	\item It provides a \emph{background} for those requirements and makes them easier to understand.
\end{itemize}


\subsection{Product Perspective}
\label{sub:product_perspective}
% Begin SubSection
\begin{itemize}
	\item Put the product into perspective with other related products, i.e., context
	\item If the product is independent and totally self-contained, it should be stated here
	\item If the SRS defines a product that is a component of a larger system, then this subsection should relate the requirements of that larger system to the functionality of the software being developed. Identify interfaces between that larger system and the software to be developed.
	\item A block diagram showing the major components of the larger system, interconnections, and external interfaces can be helpful
\end{itemize}
% End SubSection

\subsection{Product Functions}
\label{sub:product_functions}
% Begin SubSection
\begin{itemize}
	\item Provide a \emph{summary} of the major functions that the software will perform.
	\begin{itemize}
		\item \textbf{Example}: An SRS for an accounting program may use this part to address customer account maintenance, customer statement, and invoice preparation without mentioning the vast amount of detail that each of those functions requires.
	\end{itemize}
	\item Functions should be organized in a way that makes the list of functions understandable to the customer or to anyone else reading the document for the first time 
	\item Present the functions in a list format - each item should be one function, with a brief description of it
	\item Textual or graphical methods can be used to show the different functions and their relationships
	\begin{itemize}
		\item Such a diagram is not intended to show a design of a product, but simply shows the logical relationships among variables
	\end{itemize} 
\end{itemize}
% End SubSection

\subsection{User Characteristics}
\label{sub:user_characteristics}
% Begin SubSection
\begin{itemize}
	\item Describe those general characteristics of the intended users of the product including educational level, experience, and technical expertise 
	\item Since there will be many users, you may wish to divide into different user types or personas
%	\item Do not state specific requirements, but rather provide the reasons why certain specific requirements are later specified
\end{itemize}
% End SubSection

\subsection{Constraints}
\label{sub:constraints}
% Begin SubSection
\begin{itemize}
	\item Provide a general description of any constraints that will limit the developer's options
\end{itemize}
% End SubSection

\subsection{Assumptions and Dependencies}
\label{sub:assumptions_and_dependencies}
% Begin SubSection
\begin{itemize}
	\item List any assumptions you made in interpreting what the software being developed is aiming to achieve
	\item List any other assumptions you made that, if it fails to hold, could require you to change the requirements
	%\item List each of the factors that affect the requirements stated in the SRS
	%\item These factors are not design constraints on the software but are, rather, any changes to them that can affect the requirements in the SRS
	\begin{itemize}
		\item \textbf{Example}: An assumption may be that a specific operating system will be available on the hardware designated for the software product. If, in fact, the operating system is not available, the SRS would then have to change accordingly.
	\end{itemize}
\end{itemize}
% End SubSection

\subsection{Apportioning of Requirements}
\label{sub:apportioning_of_requirements}
% Begin SubSection
\begin{itemize}
	\item Identify requirements that may be delayed until future versions of the system
\end{itemize}
% End SubSection

% End Section
\section{Use Case Diagram}
\label{sec:use_case_diagram}
% Begin Section
\begin{itemize}
	\item Provide the use case diagram for the system being developed.
	\item You do not need to provide the textual description of any of the use cases here (these will be specified under "Highlights of Functional Requirements").
%	\item Provide \emph{one} use case diagram for the most important Business Event.
%	\item The text of all use cases will be specified under "Highlights of Functional Requirements"
\end{itemize}
%In this section, select the most important Business Event that your system responds to and give its use case diagram.  Only one use case diagram is needed.  Give a brief textual description of the use case without repeating what is in the scenarios of the corresponding Business Event.

%
%
%
%This section should provide a use case diagram for your application. 
%\begin{enumerate}[a)]
%	\item Each use case appearing in the diagram should be accompanied by a text description. 
%\end{enumerate}
%% End Section

\section{Highlights of Functional Requirements}
\label{sec:functional_requirements}
% Begin Section
\begin{itemize}
	\item Specify all use cases (or other scenarios triggered by other events), organized by Business Event. 
	\item For each Business Event, show the scenario from every Viewpoint. You should have the same set of Viewpoints across all Business Events. If a Viewpoint doesn't participate, write N/A so we know you considered it still. You can choose how to present this - keep in mind it should be easy to follow. 
	\item At the end, combine them all into a Global Scenario.
	%\item Specify the "use cases" (or other triggering events) organized by Business Event. (The Global Scenario is what you might think of as a use case). Be sure to consider Business Events that aren't just triggered by users with goals (e.g. something happens in the environment that your system needs to respond to)
	\item Your focus should be on what the system needs to do, not how to do it. Specify it in enough detail that it clearly specifies what needs to be accomplished, but not so detailed that you start programming or making design decisions.
	\item Keep the length of each use case (Global Scenario) manageable. If it's getting too long, split into sub-cases.
	\item You are \emph{not} specifying a complete and consistent set of functional requirements here. (i.e. you are providing them in the form of use cases/global scenarios, not a refined list). For the purpose of this project, you do not need to reduce them to a list; the global scenarios format is all you need.
	\item Red text below is just to highlight where you need to insert a scenario - don't actually write it all in red.
\end{itemize}

\noindent {\bf Main Business Events:} List out all the main business events you are presenting. If you sub-divided into smaller ones, you don't need to include the smaller ones in this list.\\

\noindent {\bf Viewpoints:} List out all the viewpoints you will be considering.\\

\noindent {\bf Interpretation:} Specify any liberties you took in interpreting business events, if necessary.\\

%%BUSINESS EVENT 1
	
	\begin{enumerate}[\bf {BE}1.]
		\item Authorizing an employee to use the chat application for work purposes.
		\item {\bf VP3: Organization}
			
			\noindent\fbox{%
				\parbox{0.5\textwidth}{%
					\begin{itemize}
						\item {\bf $S_{1}$:} The organization provides a list of authorized individuals to the IT Support team.
					\end{itemize}
				}%
			}
                \item {\bf VP2: IT Support Team}
			
			\noindent\fbox{%
				\parbox{0.5\textwidth}{%
					\begin{itemize}
						\item {\bf $S_{1}$:} The IT Support team grants authorization to the individuals to gain access and use the application.
					\end{itemize}
				}%
			}
%			
		
		\item[] {\bf Global Scenario of {\it Authorizing an employee to use the chat application for work purposes.}:} \newline
		\noindent\fbox{%
			\parbox{0.5\textwidth}{%
				\begin{itemize}
                        \item {\bf $S_{1}$:} The organization provides a list of authorized individuals to the IT Support team. \textcolor{green}{Proceed to E.1}.
					  \item {\bf $E_{1}$:} The IT Support team grants authorization to the individuals to gain access and use the application.\textcolor{green}{Proceed to S.2}.
					  \item {\bf $S_{2}$:} Employees now have access to the system and are authorized to use it.
                            
						
						
				\end{itemize}
			}%
		}%
		\end{enumerate}


%%BUSINESS EVENT 2
	
	\begin{enumerate}[\bf {BE}2.]
		\item Logging onto the chat application
		\item {\bf VP1: Employee}
			
			\noindent\fbox{%
				\parbox{0.5\textwidth}{%
					\begin{itemize}
						\item {\bf $S_{1}$:} System displays required login fields.
						\item {\bf $E_{1}$:} User inputs personal login information.
						\item {\bf $S_{2}$:}  Information is verified by the system. 
						\item {\bf $E_{2}$:}  User is sent a verification code. 
						
						\item {\bf $S_{3}$:}  User enters code. 
						\item {\bf $S_{(4)}$:} System reviews code and if correct, allows user access to the application.
					\end{itemize}
				}%
			}
%			
		
		\item[] {\bf Global Scenario of {\it Logging onto the chat application}:} \newline
		\noindent\fbox{%
			\parbox{0.5\textwidth}{%
				\begin{itemize}
                        \item {\bf $S_{1}$:} System displays required login fields. \textcolor{green}{Proceed to E.1}.
					  \item {\bf $E_{1}$:} User inputs personal login information.\textcolor{green}{Proceed to S.2}.
					  \item {\bf $S_{2}$:}  Information is verified by the system.
                            \begin{itemize}
                                {\bf $S_{2}.1$} The user information is inaccurate. \textcolor{red}{Loop to S.1}.
                                {\bf $S_{2}.2$} The user information is accurate. \textcolor{green}{Proceed to E.2}.
                                \end{itemize}
						\item {\bf $E_{2}$:}  User is sent a verification code.  \textcolor{green}{Proceed to S.3}.
						
						\item {\bf $S_{3}$:}  User enters code. 
                            \begin{itemize}
                                \item {\bf $S_{3}.1$} The entered code is wrong. \textcolor{red}{Loop to E.2}.
                                \item {\bf $S_{3}.2$} The user did not receive a code. \textcolor{red}{Proceed to E.2}.
                                \item {\bf $S_{3}.3$} The code entered is correct. \textcolor{green}{Proceed to S.4}.
                            \end{itemize}
					\item {\bf $S_{(4)}$:} User is allowed access to the application.
				\end{itemize}
			}%
		}%
		\end{enumerate}
  
  %%BUSINESS EVENT 3
  \begin{enumerate}[\bf {BE}3.]
		\item Creating individual chat threads
		\item {\bf VP1: Employee}
  
		\noindent\fbox{%
			\parbox{0.5\textwidth}{%
				\begin{itemize}
						\item {\bf $S_{1}$:} User logs into the application.
						\item {\bf $S_{2}$:}  User searches up the individual they want to create a chat with.
						\item {\bf $S_{3}$:}  User selects the create chat option and sends a chat to their colleague.
						\item {\bf $E_{(1)}$:} Colleague receives chat request. The chat thread is now created.
					\end{itemize}
				}%
			}
            \item {\bf VP2: IT Support Team}
  
			\noindent\fbox{%
				\parbox{0.5\textwidth}{%
					\begin{itemize}
						\item {\bf $S_{1}$:} IT Support Team receives a notification for an open ticket.
						\item {\bf $S_{2}$:} IT Support Team member resolves the ticket.
                            \item {\bf $S_{3}$:} The ticket is closed.
					\end{itemize}
				}%
			}
%			
		
		\item[] {\bf Global Scenario of {\it Creating chat threads with an individual}:} \newline
           
                       
		\noindent\fbox{%
			\parbox{0.5\textwidth}{%
				\begin{itemize}
					\item {\bf $S_{1}$:} User logs into the application.\textcolor{green}{Proceed to S.2}.
					\item {\bf $S_{2}$:}  User searches up the individual they want to create a chat with.
                        \begin{itemize}
                            \item {\bf $S_{2}.1$} Their colleague is not in the application. \textcolor{red}{Proceed to S.3.1}.
                            \item {\bf $S_{2}.2$} Their colleague is able to use the application and is in the system. \textcolor{green}{Proceed to E.1}.
                        \end{itemize}
                    \item {\bf $S_{3}$:}  The IT Support team is notified, and the required colleague is granted authorization to use the system.\textcolor{green}{Proceed to E.1}.
                    \item {\bf $E_{1}$:}  Colleague receives chat invite. Chat thread is created. 
                            
			\end{itemize}
			}%
		}%	
		\end{enumerate}

  %%BUSINESS EVENT 4
  \begin{enumerate}[\bf {BE}4.]
		\item Creating chat threads with multiple people
		\item {\bf VP1: Employee}
  
			\noindent\fbox{%
				\parbox{0.5\textwidth}{%
					\begin{itemize}
						\item {\bf $S_{1}$:} User logs into the application.
						\item {\bf $S_{2}$:}  User searches up the individuals they want to create a chat with.
						\item {\bf $S_{3}$:}  User selects the create chat option and sends a chat to their colleagues.
						\item {\bf $E_{(1)}$:} Colleagues receives chat request. The chat thread is now created.
					\end{itemize}
				}%
			}
            \item {\bf VP2: IT Support Team}
  
			\noindent\fbox{%
				\parbox{0.5\textwidth}{%
					\begin{itemize}
						\item {\bf $S_{1}$:} IT Support Team receives a notification for an open ticket.
						\item {\bf $S_{2}$:} IT Support Team member resolves the ticket.
                            \item {\bf $S_{3}$:} The ticket is closed.
					\end{itemize}
				}%
			}
%			
		
		\item[] {\bf Global Scenario of {\it Creating chat threads with multiple people}:} \newline
		\noindent\fbox{%
			\parbox{0.5\textwidth}{%
				\begin{itemize}
					\item {\bf $S_{1}$:} User logs into the application.\textcolor{green}{Proceed to S.2}.
					\item {\bf $S_{2}$:}  User searches up the individuals they want to create a chat with. 
                        \begin{itemize}
                            \item {\bf $S_{2}.1$} Their colleagues are not in the application. \textcolor{red}{Proceed to S.3}.
                            \item {\bf $S_{2}.2$} Their colleagues are able to use the application and is in the system. \textcolor{green}{Proceed to E.1}.
                        \end{itemize}
                    \item {\bf $S_{3}$:}  The IT Support team is notified, and the required colleagues are granted authorization to use the system.\textcolor{green}{Proceed to E.1}.
						
					\item {\bf $E_{1}$:}  Colleagues receive chat invite. Chat thread is created. 
                            
				\end{itemize}
			}%
		}%	
		\end{enumerate}
		
%%BUSINESS EVENT 5
  \begin{enumerate}[\bf {BE}5.]
		\item Sending text-only chats
		\item {\bf VP1: Employee}
  
			\noindent\fbox{%
				\parbox{0.5\textwidth}{%
					\begin{itemize}
						\item {\bf $S_{1}$:} User opens the relevant chat thread.
						\item {\bf $S_{2}$:}  User utilizes the typing box to write a text-based message to their colleague(s).
						\item {\bf $S_{3}$:}  User presses the send button to deliver the communication.
						\item {\bf $E_{(1)}$:} Colleague(s) receive the chat(s).
					\end{itemize}
				}%
			}
                \item {\bf VP2: IT Support Team}
  
			\noindent\fbox{%
				\parbox{0.5\textwidth}{%
					\begin{itemize}
						\item {\bf $S_{1}$:} IT Support Team receives a notification for an open ticket.
						\item {\bf $S_{2}$:} IT Support Team member resolves the ticket.
                            \item {\bf $S_{3}$:} The ticket is closed.
					\end{itemize}
				}%
			}
%			
		
		\item[] {\bf Global Scenario of {\it Sending text-only chats}:} \newline
		\noindent\fbox{%
			\parbox{0.5\textwidth}{%
				\begin{itemize}
					\item {\bf $S_{1}$:} User opens the relevant chat thread.\textcolor{green}{Proceed to S.2}.
					\item {\bf $S_{2}$:}   User utilizes the typing box to write a text-based message to their colleague(s).\textcolor{green}{Proceed to S.3}.
                     \item {\bf $S_{3}$:}  User presses the send button to deliver the communication.
                        \begin{itemize}
                            \item {\bf $S_{3}.1$} The chat does not send. \textcolor{red}{Proceed to S.4}.
                            \item {\bf $S_{3}.2$} The chat shows as sent, but the other members of the chat do not receive it. \textcolor{red}{Proceed to S.4}.
                            \item {\bf $S_{3}.3$} The chat sends successfully. \textcolor{green}{Proceed to E.1}.
                        \end{itemize}
                    \item {\bf $S_{4}$:}  The IT Support team is notified. \textcolor{green}{Proceed to E.1}.
					\item {\bf $E_{1}$:}  Colleague receives chat, and can respond accordingly.
                            
				\end{itemize}
			}%
		}%	
		\end{enumerate}
		
%%BUSINESS EVENT 6
  \begin{enumerate}[\bf {BE}6.]
		\item Sending file attachments in communication.
		\item {\bf VP1: Employee}
  
			\noindent\fbox{%
				\parbox{0.5\textwidth}{%
					\begin{itemize}
						\item {\bf $S_{1}$:} User opens the relevant chat thread.
						\item {\bf $S_{2}$:} User utilizes the attachment button to select the file(s) they wish to send.
						\item {\bf $S_{3}$:}  User can also drag and drop the file(s) they wish to send in the chat box.
                            \item {\bf $S_{4}$:}  User presses the send button to deliver the communication.
						\item {\bf $E_{(1)}$:} Colleague(s) receive the file attachment(s).
					\end{itemize}
				}%
			}
                \item {\bf VP2: IT Support Team}
  
			\noindent\fbox{%
				\parbox{0.5\textwidth}{%
					\begin{itemize}
						\item {\bf $S_{1}$:} IT Support Team receives a notification for an open ticket.
						\item {\bf $S_{2}$:} IT Support Team member resolves the ticket.
                            \item {\bf $S_{3}$:} The ticket is closed.
					\end{itemize}
				}%
			}
%			
		
		\item[] {\bf Global Scenario of {\it Sending file attachments in communication.}:} \newline
		\noindent\fbox{%
			\parbox{0.5\textwidth}{%
				\begin{itemize}
					\item {\bf $S_{1}$:} User opens the relevant chat thread.\textcolor{green}{Proceed to S.2}.
					\item {\bf $S_{2}$:}  User utilizes the attachment button to select the file(s) they wish to send. \textcolor{green}{Proceed to S.3}.
                    \item {\bf $S_{3}$:} User can also drag and drop the file(s) from their laptop they wish to send in the chat box.\textcolor{green}{Proceed to S.4}.
                    \item {\bf $S_{4}$:}  User presses the send button to deliver the communication.
                        \begin{itemize}
                            \item {\bf $S_{4}.1$} The file does not send. \textcolor{red}{Proceed to E.1}.
                            \item {\bf $S_{4}.2$} The file shows as sent, but the other members of the chat do not receive it. \textcolor{red}{Proceed to E.1}.
                            \item {\bf $S_{4}.2$} The file is sent and received; however, the receiving individuals are unable to open and view the file. \textcolor{red}{Proceed to E.1}.
                             \item {\bf $S_{4}.3$} The file sends successfully. \textcolor{green}{Proceed to E.2}.
                        \end{itemize}
                    \item {\bf $E_{1}$:}  A ticket is opened with the IT Support team to resolve the issue.
                        \begin{itemize}
                        
                        \item {\bf $E_{1}.1$} The issue is not resolved. \textcolor{red}{Loop to E.1}.
                        \item {\bf $E_{1}.2$} The issue is resolved and the file sends successfully. \textcolor{green}{Proceed to E.2}.
                        \end{itemize}
                    \item {\bf $E_{2}$:}  Colleague receives file, and can open and view the file.
				\end{itemize}
			}%
		}%	
		\end{enumerate}

%%BUSINESS EVENT 7
  \begin{enumerate}[\bf {BE}7.]
		\item Adding individuals to a pre-existing chat
		\item {\bf VP1: Employee}
  
			\noindent\fbox{%
				\parbox{0.5\textwidth}{%
					\begin{itemize}
						\item {\bf $S_{1}$:} User opens the relevant chat thread.
						\item {\bf $S_{2}$:} User searches up the individuals they want to add to the chat.
						\item {\bf $S_{3}$:}  User selects the add individual option to add them to the chat.
					\end{itemize}
				}%
			}

                \item {\bf VP2: IT Support Team}
  
			\noindent\fbox{%
				\parbox{0.5\textwidth}{%
					\begin{itemize}
						\item {\bf $S_{1}$:} IT Support Team receives a notification for an open ticket.
						\item {\bf $S_{2}$:} IT Support Team member resolves the ticket.
                            \item {\bf $S_{3}$:} The ticket is closed.
					\end{itemize}
				}%
			}
%			
		
		\item[] {\bf Global Scenario of {\it Adding individuals to a pre-existing group chat.}:} \newline
		\noindent\fbox{%
			\parbox{0.5\textwidth}{%
				\begin{itemize}
					\item {\bf $S_{1}$:} User opens the relevant chat thread.\textcolor{green}{Proceed to S.2}.
					\item {\bf $S_{2}$:}   User searches up the individuals they want to add to the chat.
                        \begin{itemize}
                            \item {\bf $S_{2}.1$:} The individual is not found. \textcolor{red}{Proceed to E.1}.
                            \item {\bf $S_{2}.2$:} The individual is in the system. \textcolor{green}{Proceed to S.3}.
                        \end{itemize}
                    \item {\bf $S_{3}$:}  User selects the add individual option to add them to the chat.
                        \begin{itemize}
                            \item {\bf $S_{3}.1$:} The individual cannot be added to the chat. \textcolor{red}{Proceed to E.1}.
                            \item {\bf $S_{3}.2$:} The individual is added to the system. \textcolor{green}{Proceed to E.2}.
                        \end{itemize}
                            
                    \item {\bf $E_{1}$:}  A ticket is opened with the IT Support team to resolve the issue.
                     \begin{itemize}
                        \item {\bf $E_{1}.1$} The issue is not resolved. \textcolor{red}{Proceed to E.1}.
                        \item {\bf $E_{1}.2$} The issue is resolved and the colleague can be added to the chat successfully. \textcolor{green}{Proceed to E.2}.
                    \end{itemize}
		\item {\bf $E_{2}$:}  Colleague is added to the chat, and can receive and send messages.
                            
				\end{itemize}
			}%
		}%	
		\end{enumerate}

%%BUSINESS EVENT 8
  \begin{enumerate}[\bf {BE}8.]
		\item Removing individuals from a pre-existing chat
		\item {\bf VP1: Employee}
  
			\noindent\fbox{%
				\parbox{0.5\textwidth}{%
					\begin{itemize}
						\item {\bf $S_{1}$:} User opens the relevant chat thread.
						\item {\bf $S_{2}$:} User searches up the individuals they want to remove from the chat.
						\item {\bf $S_{3}$:}  User selects the remove individual option to remove them from the chat.
					\end{itemize}
				}%
			}
                \item {\bf VP2: IT Support Team}
  
			\noindent\fbox{%
				\parbox{0.5\textwidth}{%
					\begin{itemize}
						\item {\bf $S_{1}$:} IT Support Team receives a notification for an open ticket.
						\item {\bf $S_{2}$:} IT Support Team member resolves the ticket.
                            \item {\bf $S_{3}$:} The ticket is closed.
					\end{itemize}
				}%
			}
%			
		
		\item[] {\bf Global Scenario of {\it Removing individuals to a pre-existing group chat.}:} \newline
		\noindent\fbox{%
			\parbox{0.5\textwidth}{%
				\begin{itemize}
					\item {\bf $S_{1}$:} User opens the relevant chat thread. \textcolor{green}{Proceed to S.1}.
					\item {\bf $S_{2}$:}   User searches up the individuals they want to remove from the chat.
                        \begin{itemize}
                            \item {\bf $S_{2}.1$} The individual is not found. \textcolor{red}{Proceed to E.1}.
                            \item {\bf $S_{2}.2$} The individual is in the chat. \textcolor{green}{Proceed to S.3}.
                        \end{itemize}
                    \item {\bf $S_{3}$:}  User selects the remove individual option to remove them from the chat.
                        \begin{itemize}
                            \item {\bf $S_{3}.1$} The individual cannot be removed from the chat. \textcolor{red}{Proceed to E.1}.
                            \item {\bf $S_{3}.2$} The individual is removed from the chat. \textcolor{green}{Proceed to E.2}.
                        \end{itemize}
                            
                    \item {\bf $E_{1}$:}  A ticket is opened with the IT Support team to resolve the issue.
                        \begin{itemize}
                            \item {\bf $E_{1}.1$} The issue is not resolved. \textcolor{red}{Loop to E.1}.
                            \item {\bf $E_{1}.2$} The issue is resolved and the colleague can be removed from the chat successfully. \textcolor{green}{Proceed to E.2}.
                        \end{itemize}
                        
					\item {\bf $E_{2}$:}  Colleague is removed from the chat, and can no longer receive or send messages, or see/have access to prior communication in that chat thread.
                            
			\end{itemize}
			}%
		}%	
		\end{enumerate}

%%BUSINESS EVENT 9
  \begin{enumerate}[\bf {BE}9.]
		\item Unsending a chat/communication
		\item {\bf VP1: Employee}
  
			\noindent\fbox{%
				\parbox{0.5\textwidth}{%
					\begin{itemize}
						\item {\bf $S_{1}$:} User opens the relevant chat with the message(s) they want to unsend.
						\item {\bf $S_{2}$:} User finds the message(s) they want to unsend.
						\item {\bf $S_{3}$:} User double clicks on the message, and selects the unsend option.
                            \item {\bf $S_{4}$:} The message disappears and is replaced with a ‘You have unsent this message’ notification.
					\end{itemize}
				}%
			}

                \item {\bf VP2: IT Support Team}
  
			\noindent\fbox{%
				\parbox{0.5\textwidth}{%
					\begin{itemize}
						\item {\bf $S_{1}$:} IT Support Team receives a notification for an open ticket.
						\item {\bf $S_{2}$:} IT Support Team member resolves the ticket.
                            \item {\bf $S_{3}$:} The ticket is closed.
					\end{itemize}
				}%
			}
%			
		
		\item[] {\bf Global Scenario of {\it Unsending a chat/communication}:} \newline
		\noindent\fbox{%
			\parbox{0.5\textwidth}{%
				\begin{itemize}
					   \item {\bf $S_{1}$:} User opens the relevant chat with the message(s) they want to unsend.\textcolor{green}{Proceed to S.2}.
					   \item {\bf $S_{2}$:} User finds the message(s) they want to unsend. \textcolor{green}{Proceed to S.3}.
                        \item {\bf $S_{3}$:} User double clicks on the message, and selects the unsend option.
                            \begin{itemize}
                                \item {\bf $S_{3}.1$} The message is not unsent. \textcolor{red}{Proceed to E.1}.
                                \item {\bf $S_{3}.2$} The message is successfully unsent. \textcolor{green}{Proceed to E.2}.
                            \end{itemize}
                            
                        \item {\bf $E_{1}$:}  A ticket is opened with the IT Support team to resolve the issue.
                            \begin{itemize}
                                \item {\bf $E_{1}.1$} The issue is not resolved. \textcolor{red}{Proceed to E.1}.
                                \item {\bf $E_{1}.2$} The issue is resolved and the colleague can be removed from the chat successfully. \textcolor{green}{Proceed to E.2}.
                            \end{itemize}
                        
					   \item {\bf $E_{2}$:}  The message is unsent, and the other members of the chat can no longer view it.
                            
				\end{itemize}
			}%
		}%	
		\end{enumerate}

%%BUSINESS EVENT 10
  \begin{enumerate}[\bf {BE}10.]
		\item Signing out of the chat application
		\item {\bf VP1: Employee}
  
			\noindent\fbox{%
				\parbox{0.5\textwidth}{%
					\begin{itemize}
						\item {\bf $S_{1}$:} User goes to their profile.
						\item {\bf $E_{1}$:} User selects the logout option.
						\item {\bf $S_{(4)}$:} User is successfully logged out.
					\end{itemize}
				}%
			}

                \item {\bf VP2: IT Support Team}
  
			\noindent\fbox{%
				\parbox{0.5\textwidth}{%
					\begin{itemize}
						\item {\bf $S_{1}$:} IT Support Team receives a notification for an open ticket.
						\item {\bf $S_{2}$:} IT Support Team member resolves the ticket.
                            \item {\bf $S_{3}$:} The ticket is closed.
					\end{itemize}
				}%
			}
%			
		
		\item[] {\bf Global Scenario of {\it Signing out of the chat application}:} \newline
		\noindent\fbox{%
			\parbox{0.5\textwidth}{%
				\begin{itemize}
					   \item {\bf $S_{1}$:} User goes to their profile. \textcolor{green}{Proceed to E.1}.
					   \item {\bf $E_{1}$:} User selects the logout option.
                            \begin{itemize}
                                \item {\bf $E_{1}.1$} The user is unable to log out. \textcolor{red}{Proceed to E.1}.
                                \item {\bf $E_{1}.2$} The user successfully logs out. \textcolor{green}{Proceed to E.3}.
                            \end{itemize}
                        \item {\bf $E_{2}$:}  A ticket is opened with the IT Support team to resolve the issue.
                        \begin{itemize}
                            \item {\bf $E_{2}.1$} The issue is not resolved. \textcolor{red}{Loop to E.2}.
                            \item {\bf $E_{2}.2$} The issue is resolved and the colleague can be removed from the chat successfully. \textcolor{green}{Proceed to E.3}.
                        \end{itemize}
					  \item {\bf $E_{3}$:}  The message is unsent, and the other members of the chat can no longer view it.
                            
					\item {\bf $E_{4}$:}  The user is shown a message saying 'You have been logged out' along with a link to go back to the sign in page. $S_{2}$
     \end{itemize}
						
				
			}%
		}%
		\end{enumerate}

%%BUSINESS EVENT 11
  \begin{enumerate}[\bf {BE}11.]
		\item Notifications being sent to users.
		\item {\bf VP1: Employee}
  
			\noindent\fbox{%
				\parbox{0.5\textwidth}{%
					\begin{itemize}
						\item {\bf $S_{1}$:} User receives a notification when a chat is directed to them, or if they are mentioned in a chat.
						\item {\bf $S_{2}$:} User clicks on the notification and is directed to the chat in question.
					\end{itemize}
				}%
			}

                \item {\bf VP2: IT Support Team}
  
			\noindent\fbox{%
				\parbox{0.5\textwidth}{%
					\begin{itemize}
						\item {\bf $S_{1}$:} IT Support Team receives a notification for an open ticket.
						\item {\bf $S_{2}$:} IT Support Team member resolves the ticket.
                            \item {\bf $S_{3}$:} The ticket is closed.
					\end{itemize}
				}%
			}
			
		
		\item[] {\bf Global Scenario of {\it Signing out of the chat application}:} \newline
		\noindent\fbox{%
			\parbox{0.5\textwidth}{%
				\begin{itemize}
					   \item {\bf $S_{1}$:} User receives a notification when a chat is directed to them, or if they are mentioned in a chat.\textcolor{green}{Proceed to S.2}.
					   \item {\bf $S_{2}$:} User clicks on the notification and is directed to the chat in question.
                            \begin{itemize}
                                \item {\bf $S_{2}.1$} The user cannot click on the notification to go to the relevant chat. \textcolor{red}{Proceed to E.2}.
                                \item {\bf $S_{2}.2$} The user can view the relevant chat after clicking on the notification. \textcolor{green}{Proceed to E.3}.
                            \end{itemize}

                    \item {\bf $E_{2}$:}  A ticket is opened with the IT Support team to resolve the issue.
                    \begin{itemize}
                        \item {\bf $E_{2}.1$} The issue is not resolved. \textcolor{red}{Loop to E.2}.
                        \item {\bf $E_{2}.2$} The issue is resolved and the colleague can be removed from the chat successfully. \textcolor{green}{Proceed to E.3}.
                    \end{itemize}
					\item {\bf $E_{3}$:}  The user is directed to the chat, and can respond accordingly. The notification banner disappears after the user has addressed it.
			\end{itemize}
				
			}%
		}%
		\end{enumerate}
		
%End Section

\section{Non-Functional Requirements}
\label{sec:non-functional_requirements}


\begin{itemize}
	\item For each non-functional requirement, provide a justification/rationale for it.\\
	{\bf Example:} \\
	SC1. \emph{The device should not explode in a customer’s pocket.}\\
	{\bf Rationale:} Other companies have had issues with the batteries they used in their phones randomly exploding [insert citation]. This causes a safety issue, as the phone is often carried in a person's hand or pocket.	
	\item If you need to make a guess because you couldn't really talk to stakeholders, you can say "We imagined stakeholders would want...because..."
	\item Each requirement should have a unique label/number for it.
	\item In the list below, if a particular section doesn't apply, just write N/A so we know you considered it.
\end{itemize}

% Begin Section
\subsection{Look and Feel Requirements}
\label{sub:look_and_feel_requirements}
% Begin SubSection

\subsubsection{Appearance Requirements}
\label{ssub:appearance_requirements}
% Begin SubSubSection
\begin{enumerate}[{LF-A}1. ]
	\item 
\end{enumerate}
% End SubSubSection

\subsubsection{Style Requirements}
\label{ssub:style_requirements}
% Begin SubSubSection
\begin{enumerate}[{LF-S}1. ]
	\item 
\end{enumerate}
% End SubSubSection

% End SubSection

\subsection{Usability and Humanity Requirements}
\label{sub:usability_and_humanity_requirements}
% Begin SubSection

\subsubsection{Ease of Use Requirements}
\label{ssub:ease_of_use_requirements}
% Begin SubSubSection
\begin{enumerate}[{UH-EOU}1. ]
	\item 
\end{enumerate}
% End SubSubSection

\subsubsection{Personalization and Internationalization Requirements}
\label{ssub:personalization_and_internationalization_requirements}
% Begin SubSubSection
\begin{enumerate}[{UH-PI}1. ]
	\item 
\end{enumerate}
% End SubSubSection

\subsubsection{Learning Requirements}
\label{ssub:learning_requirements}
% Begin SubSubSection
\begin{enumerate}[{UH-L}1. ]
	\item 
\end{enumerate}
% End SubSubSection

\subsubsection{Understandability and Politeness Requirements}
\label{ssub:understandability_and_politeness_requirements}
% Begin SubSubSection
\begin{enumerate}[{UH-UP}1. ]
	\item 
\end{enumerate}
% End SubSubSection

\subsubsection{Accessibility Requirements}
\label{ssub:accessibility_requirements}
% Begin SubSubSection
\begin{enumerate}[{UH-A}1. ]
	\item 
\end{enumerate}
% End SubSubSection

% End SubSection

\subsection{Performance Requirements}
\label{sub:performance_requirements}
% Begin SubSection

\subsubsection{Speed and Latency Requirements}
\label{ssub:speed_and_latency_requirements}
% Begin SubSubSection
\begin{enumerate}[{PR-SL}1. ]
	\item 
\end{enumerate}
% End SubSubSection

\subsubsection{Safety-Critical Requirements}
\label{ssub:safety_critical_requirements}
% Begin SubSubSection
\begin{enumerate}[{PR-SC}1. ]
	\item 
\end{enumerate}
% End SubSubSection

\subsubsection{Precision or Accuracy Requirements}
\label{ssub:precision_or_accuracy_requirements}
% Begin SubSubSection
\begin{enumerate}[{PR-PA}1. ]
	\item 
\end{enumerate}
% End SubSubSection

\subsubsection{Reliability and Availability Requirements}
\label{ssub:reliability_and_availability_requirements}
% Begin SubSubSection
\begin{enumerate}[{PR-RA}1. ]
	\item 
\end{enumerate}
% End SubSubSection

\subsubsection{Robustness or Fault-Tolerance Requirements}
\label{ssub:robustness_or_fault_tolerance_requirements}
% Begin SubSubSection
\begin{enumerate}[{PR-RFT}1. ]
	\item 
\end{enumerate}
% End SubSubSection

\subsubsection{Capacity Requirements}
\label{ssub:capacity_requirements}
% Begin SubSubSection
\begin{enumerate}[{PR-C}1. ]
	\item 
\end{enumerate}
% End SubSubSection

\subsubsection{Scalability or Extensibility Requirements}
\label{ssub:scalability_or_extensibility_requirements}
% Begin SubSubSection
\begin{enumerate}[{PR-SE}1. ]
	\item 
\end{enumerate}
% End SubSubSection

\subsubsection{Longevity Requirements}
\label{ssub:longevity_requirements}
% Begin SubSubSection
\begin{enumerate}[{PR-L}1. ]
	\item 
\end{enumerate}
% End SubSubSection

% End SubSection

\subsection{Operational and Environmental Requirements}
\label{sub:operational_and_environmental_requirements}
% Begin SubSection

\subsubsection{Expected Physical Environment}
\label{ssub:expected_physical_environment}
% Begin SubSubSection
\begin{enumerate}[{OE-EPE}1. ]
	\item 
\end{enumerate}
% End SubSubSection

\subsubsection{Requirements for Interfacing with Adjacent Systems}
\label{ssub:requirements_for_interfacing_with_adjacent_systems}
% Begin SubSubSection
\begin{enumerate}[{OE-IA}1. ]
	\item 
\end{enumerate}
% End SubSubSection

\subsubsection{Productization Requirements}
\label{ssub:productization_requirements}
% Begin SubSubSection
\begin{enumerate}[{OE-P}1. ]
	\item 
\end{enumerate}
% End SubSubSection

\subsubsection{Release Requirements}
\label{ssub:release_requirements}
% Begin SubSubSection
\begin{enumerate}[{OE-R}1. ]
	\item 
\end{enumerate}
% End SubSubSection

% End SubSection

\subsection{Maintainability and Support Requirements}
\label{sub:maintainability_and_support_requirements}
% Begin SubSection

\subsubsection{Maintenance Requirements}
\label{ssub:maintenance_requirements}
% Begin SubSubSection
\begin{enumerate}[{MS-M}1. ]
	\item 
\end{enumerate}
% End SubSubSection

\subsubsection{Supportability Requirements}
\label{ssub:supportability_requirements}
% Begin SubSubSection
\begin{enumerate}[{MS-S}1. ]
	\item 
\end{enumerate}
% End SubSubSection

\subsubsection{Adaptability Requirements}
\label{ssub:adaptability_requirements}
% Begin SubSubSection
\begin{enumerate}[{MS-A}1. ]
	\item 
\end{enumerate}
% End SubSubSection

% End SubSection

\subsection{Security Requirements}
\label{sub:security_requirements}
% Begin SubSection

\subsubsection{Access Requirements}
\label{ssub:access_requirements}
% Begin SubSubSection
\begin{enumerate}[{SR-AC}1. ]
	\item 
\end{enumerate}
% End SubSubSection

\subsubsection{Integrity Requirements}
\label{ssub:integrity_requirements}
% Begin SubSubSection
\begin{enumerate}[{SR-INT}1. ]
	\item 
\end{enumerate}
% End SubSubSection

\subsubsection{Privacy Requirements}
\label{ssub:privacy_requirements}
% Begin SubSubSection
\begin{enumerate}[{SR-P}1. ]
	\item 
\end{enumerate}
% End SubSubSection

\subsubsection{Audit Requirements}
\label{ssub:audit_requirements}
% Begin SubSubSection
\begin{enumerate}[{SR-AU}1. ]
	\item 
\end{enumerate}
% End SubSubSection

\subsubsection{Immunity Requirements}
\label{ssub:immunity_requirements}
% Begin SubSubSection
\begin{enumerate}[{SR-IM}1. ]
	\item 
\end{enumerate}
% End SubSubSection

% End SubSection

\subsection{Cultural and Political Requirements}
\label{sub:cultural_and_political_requirements}
% Begin SubSection

\subsubsection{Cultural Requirements}
\label{ssub:cultural_requirements}
% Begin SubSubSection
\begin{enumerate}[{CP-C}1. ]
	\item 
\end{enumerate}
% End SubSubSection

\subsubsection{Political Requirements}
\label{ssub:political_requirements}
% Begin SubSubSection
\begin{enumerate}[{CP-P}1. ]
	\item 
\end{enumerate}
% End SubSubSection

% End SubSection

\subsection{Legal Requirements}
\label{sub:legal_requirements}
% Begin SubSection

\subsubsection{Compliance Requirements}
\label{ssub:compliance_requirements}
% Begin SubSubSection
\begin{enumerate}[{LR-COMP}1. ]
	\item 
\end{enumerate}
% End SubSubSection

\subsubsection{Standards Requirements}
\label{ssub:standards_requirements}
% Begin SubSubSection
\begin{enumerate}[{LR-STD}1. ]
	\item 
\end{enumerate}
% End SubSubSection

% End SubSection

% End Section

\section{Innovative Features}
\label{sec:innovative_features}
%Begin Section

A list of potential innovative features that our group came up with has been provided below:
\begin{itemize}
	\item File-sharing for audio files (such as voice notes), image files, PDFs, emojis and GIFs.
 	\begin{itemize}
   		\item Customizable permissions for shared files (view only, edit etc.)
     	\end{itemize}
     	\item Customizable user profiles including preferred name, role in the organization and custom status.
      	\item Creating group chats with up to 50 users.
       	\begin{itemize}
		\item Customizable group details including group name and group description.
  		\item Customizable group permissions such as admin-only edit access, admin-only messaging etc.
    	\end{itemize}
     	\item An Outlook and Gmail plug-in that allows users to embed an email or email thread to a chat for easier accesss.
      	\item Automatic summarization feature for group chats that summarizes the messages received based on topics of conversation along with a short snippet and a list of users involved.
	\item Deleting and editing messages within 10 minutes of sending it.
 	\item In-built biometric authentication feature for general app as well as individual chats to ensure increased security.
  	\begin{itemize}
   		\item Back-up pin protection in cases of biometric failure.
     	\end{itemize}
  	\item Message broadcasting that allows sending a message to a group of users without creating a group chat.
   	\item An out-of-office automatic reply option that allows users to set-up a custom message to sent when they are away.
    	\item Live translation option for messages sent or received in a language other than English.
     	\item Geofencing that restricts a user's ability to send or received messages in a group chat if they are outside the permitted virtual boundary.
      	\item Smart reply suggestions for general messages.
       	\item A vanish mode option that allows users to have a conversation with self-deleting messages that disappear when the user closes the chat.
	\end{itemize}
       	
%End Section

\appendix
\section{Division of Labour}
\label{sec:division_of_labour}
% Begin Section
Include a Division of Labour sheet which indicates the contributions of each team member. This sheet must be signed by all team members.
\subsection{Awurama Nyarko}
\label{subsec:awurama_nyarko}
\begin{itemize}
	\item 1.1 Purpose
 	\item 1.2 Scope
  	\item 1.3 Definitions, Acronyms and Abbreviations
   	\item 1.4 References
    	\item 1.5 Overview
     	\item Brainstormed innovative features as a group.
\end{itemize}

\subsection{Chelsea Maramot}
\label{subsec:chelsea_maramot}
\begin{itemize}
	\item 5.1 Look and Feel Requirements
 	\item 5.2 Usability and Humanity Requirements
  	\item 5.3 Performance Requirements
   	\item 5.4 Operational and Environmental Requirements
    	\item 5.5 Maintainability and Support Requirements
     	\item 5.6 Security Requirements
      	\item 5.7 Cultural and Political Requirements
       	\item 5.8 Legal Requirements
	\item Brainstormed innovative features as a group.
	
	\includegraphics[width=0.8\textwidth]{chelsea.png}
\end{itemize}



\subsection{Harrison Chiu}
\label{subsec:harrison_chiu}
\begin{itemize}
 	\item Figure 3. Use Case Diagram
   	\item Brainstormed innovative features as a group.
 \end{itemize}

\subsection{Khushi Bhojane}
\label{subsec:khushi_bhojane}
\begin{itemize}
	\item Highlights of Functional Requirements.
 	\item Brainstormed innovative features as a group.
\end{itemize}

\subsection{Sumanya Gulati}
\label{subsec:sumanya_gulati}
\begin{itemize}
	\item 2.1 Product Perspective
 	\item Figure 1. System Diagram
  	\item 2.2 Product Functions
  	\item Figure 2. State Diagram
   	\item 2.3 User Characteristics
    	\item 2.4 Constraints
     	\item 2.5 Assumptions and Dependencies
      	\item 2.6 Apportioning of Requirements
       	\item Brainstormed innovative features as a group.
\end{itemize}
% End Section

%\newpage
%\section*{IMPORTANT NOTES}
%\begin{itemize}
%	\item Be sure to include all sections of the template in your document regardless whether you have something to write for each or not
%	\begin{itemize}
%		\item If you do not have anything to write in a section, indicate this by the \emph{N/A}, \emph{void}, \emph{none}, etc.
%	\end{itemize}
%	\item Uniquely number each of your requirements for easy identification and cross-referencing
%	\item Highlight terms that are defined in Section~1.3 (\textbf{Definitions, Acronyms, and Abbreviations}) with \textbf{bold}, \emph{italic} or \underline{underline}
%	\item For Deliverable 1, please highlight, in some fashion, all (you may have more than one) creative and innovative features. Your creative and innovative features will generally be described in Section~2.2 (\textbf{Product Functions}), but it will depend on the type of creative or innovative features you are including.
%\end{itemize}


\end{document}
%------------------------------------------------------------------------------
